\section{Marco teórico}

Sea $\Omega$ un espacio de probabilidad

\subsection{Variables aleatorias}
En estadística para una función aleatoria $F$ se hace las
siguientes definiciones:

\subsubsection{Media}
\begin{equation}
    \label{eq:media}
    \mean{F} = \int_{\Omega}F(\bullet)p(\bullet)d(\bullet) 
\end{equation}

\subsubsection{n-esimo momento central}

\begin{equation}
    \label{eq:n-momento}
    \mean{\left(F^n - \mean{F}\right)^n}
\end{equation}

\subsubsection{Desviación estandar}
\begin{equation}
    \label{eq:std}
    \sigma^2 = \mean{(F - \mean{F})^2} = \mean{F^2} - \mean{F}^2 
\end{equation}

\subsection{Errores de medición}
Supongamos que se quiere calcular el error de una medición
$E \pm a, a > 0; E = f(\bullet)$

Si se tiene una distribución uniforme ($P(\bullet) = 1)$, entonces

\begin{equation*}
    \mean{E} = \int_{-a}^{a}E P(\bullet) d(\bullet) = 
    \left[
    \frac{1}{2} E^2
    \right]_{-a}^{a} = 0
\end{equation*}

Ahora supongamos que se tiene $\bullet := (x_1, \cdots, x_n)$ y $E = f(\bullet)$,
entonces podemos aproximar $E$ mediante series de Taylor, haciendo
$\overline{\bullet} = (\overline{x_1}, \cdots, \overline{x_n})$, al rededor
de $\mean{E}$:

\newcommand{\pf}[1]{
\left[
 \frac{\partial f}{\partial x_#1}       
    \right]_{x_#1 = \overline{x}_#1}
}

\begin{equation*}
    E = \mean{E} + \sum_{i=1}^{N}
    \pf{i}(x - \overline{x}_i) + O(\bullet)
\end{equation*}

\begin{equation}
    \label{eq:aprox}
    E = \sum_{i=1}^{N}
    \pf{i}
    (x - \overline{x}_i)
\end{equation}

Por lo que la desviación estándar de $E$ es:
\begin{align*}
  u_c^2 (E) &= \sigma^2 = \mean{E^2}-\mean{E}^2 = \int_{\Omega} {\left(
    \sum_{i=1}^{N}
    \pf{i}(x - \overline{x}_i)
\right)}^2
P(\bullet)d(\bullet)\\
&= \int_\Omega \left[
    \sum_{i=1}^{N}
        \pf{i}^2{(x - \overline{x}_i)}^2
        \right]P(\bullet)d(\bullet)
        \\
        +
        &\int_{\Omega} \left[
            \sum_{i=1}^{N} \sum_{j\neq i}^{N}
        \pf{i} (x - \overline{x}_i)
        \pf{j}(x - \overline{x}_j)
        \right]
        P(\bullet)d(\bullet)
        \\
\end{align*}

Suponiendo que las probabilidades combinadas son independientes
entonces el ultimo sumando de la ecuación anterior se reduce a:

\begin{equation*}
u_c^2(E) = \int_\Omega \sum_{i=1}^{N}
        \left(
        \pf{i}
        \right)^2
        {(x - \overline{x}_i)}^2
        P(\bullet)d(\bullet)
\end{equation*}